\section{Hypotheses}

The THUMS hypotheses H1--H5 are retained. H6 is withdrawn and replaced. The
hypotheses specific to BOREAS are labelled B1--B6.

\begin{description}[leftmargin=2.2em,style=nextline]
\item[B1 --- Parallel downward tubes.] The $\nft$ finned tubes are
hydraulically in parallel, carry equal mass flow, and are thermally identical
at a given depth. Flow in them is downward on charge and upward on discharge.
\item[B2 --- Non-exchanging downcomer.] The downcomer exchanges heat with the
surrounding material only through its insulation, quasi-steadily, and does not
participate in the enthalpy march. Its leak is accounted separately
(Section~\ref{sec:bypass}).
\item[B3 --- One-dimensional formation path.] Heat leaves the borehole
radially. Axial conduction in the formation is neglected, which is justified
by the aspect ratio: the thermal penetration depth over the evaluation horizon
is \SI{46}{\metre} against a well length of \SI{1524}{\metre}.
\item[B4 --- Quasi-steady ground.] Within one charge--discharge cycle the
formation presents a fixed resistance to a fixed far-field temperature. The
cycle time is \SI{10}{\hour} and the formation time constant is years, a
separation of four orders of magnitude.
\item[B5 --- Undisturbed far field.] The reference temperature for the loss is
the \emph{undisturbed} geothermal profile. The warming of the rock immediately
around the field is represented through the time dependence of $\Rg$ and not
as a separate capacitance.
\item[B6 --- Uniform cluster.] All $N$ wells in a field are identical and
equally loaded, so the perimeter loss divides equally among them. Edge wells
in a real cluster lose more than interior ones; this is an averaging, not a
bound.
\end{description}

